% ---------------------------------------------------------------------
% ---------------------------------------------------------------------
\chapter{Du jugement à l'outil~: le protocole, ses limites, ses suites}
\label{chap:protocole}

\section{Le protocole récapitulé}
\label{sec:protocole}

Rappelons que ce travail ne fournit pas une stratégie gagnante, mais une
procédure de jugement, réutilisable telle quelle sur tout signal futur de
l'agent, qu'il soit technique, informationnel ou relationnel. Chaque étage
répare une condition d'application précise, et le verdict n'est prononcé que
si tous les étages concordent. Le tableau~\ref{tab:protocole} récapitule cette
procédure en six étages.

\begin{table}[H]
\centering
\caption{Le protocole de validation en six étages~: chaque étage répare
une menace identifiée, et le verdict exige de tous les franchir.}
\label{tab:protocole}
\small
\begin{tabular}{@{}p{0.30\linewidth}p{0.34\linewidth}p{0.26\linewidth}@{}}
\toprule
Étage & Menace traitée & Outil \\
\midrule
1. \edgevar{} apparié au hasard & dérive du marché & simulation appariée \\
2. Student unilatéral et théorème central limite & variabilité d'échantillonnage & test de référence \\
3. Seuil de Bonferroni & multiplicité (11 tests) & contrôle du risque global \\
4. Effet de sondage et rééchantillonnage par grappes & dépendance intra-titre & théorie des sondages \\
5. Les deux portes (titre, mois) & poids inégal des grappes & agrégation équipondérée \\
6. Série mensuelle autorégressive & autocorrélation résiduelle & variance de long terme \\
\bottomrule
\end{tabular}
\end{table}


\section{Limites assumées}
\label{sec:limites}

Aucune de ces limites n'invalide le verdict rendu~; chacune trace en revanche
la frontière de ce que le protocole peut affirmer. On les regroupe en limites
techniques, propres à la construction des données, et en limites statistiques,
propres aux méthodes elles-mêmes.

\subsection*{Limites techniques}

La première tient à la façon dont les stratégies ont été établies. Si l'on
étudiait des règles optimisées sur les données mêmes qui servent ensuite à les
juger, on tomberait dans le sur-apprentissage. Le risque est ici atténué, car
les stratégies retenues sont des règles classiques à peu de paramètres, non
ajustées sur l'échantillon. Il n'est pas pour autant nul, et une validation
hors échantillon, consistant à figer les règles sur une période puis à juger
sur la suivante, resterait nécessaire avant tout usage opérationnel.

Par ailleurs, les frais et les conditions d'exécution n'ont été modélisés que
de façon simplifiée. Les intégrer plus finement rapprocherait le test des
conditions réelles, et cela pèserait d'autant plus que les avantages mesurés
sont faibles.

\subsection*{Limites tenant à l'univers retenu}

Le choix du S\&P~500 comme terrain d'application appelle plusieurs réserves,
qui tiennent à la nature même de cet indice.

La première est le biais du survivant. La base ne contient que les sociétés
présentes dans l'indice au moment de sa constitution, si bien que les
entreprises ayant fait faillite ou ayant été retirées de la cote en sont
absentes. Or ce sont précisément celles dont les cours se sont effondrés. Le
biais est atténué par la construction appariée de l'\edgevar{}, qui compare
chaque opération à un tirage au hasard sur le même titre et subit donc le même
filtre, mais il subsiste.

La deuxième tient à la composition figée de l'indice. Nous appliquons aux
titres qui composent aujourd'hui le S\&P~500 des données remontant à 2010,
alors que sa composition a changé sur la période. Une société entrée dans
l'indice récemment y est entrée parce qu'elle a réussi, et lui appliquer
rétroactivement une stratégie revient à la juger sur une trajectoire dont on
sait déjà qu'elle fut favorable.

La troisième est la dérive haussière structurelle. Cet indice rassemble les
plus grandes capitalisations américaines, sélectionnées pour l'essentiel parce
qu'elles ont prospéré, ce qui explique la tendance à croître observée sur
l'ensemble de la période. C'est cette dérive qui a rendu nécessaire la
construction de l'\edgevar{} dès le premier chapitre, mais elle rend aussi le
verdict plus sévère, puisqu'une stratégie doit faire mieux qu'un marché
lui-même porteur.

La quatrième, plus anecdotique, tient au décompte des titres. L'indice compte
501 lignes de cotation pour 500 sociétés, quelques entreprises disposant de
plusieurs catégories d'actions cotées séparément. Deux lignes d'une même
société ne sont évidemment pas indépendantes, ce qui constitue une grappe que
notre découpage par titre ne capture pas.

\subsection*{Limites statistiques}

Le test de Student initial concède beaucoup d'incertitude, puisqu'il surestime
la quantité d'information disponible en traitant les opérations comme
indépendantes. Il a néanmoins suffi à écarter huit stratégies. Il suppose
aussi un effectif assez grand pour invoquer le théorème central limite,
condition d'autant plus importante qu'en finance les données ne suivent guère
une loi normale.

Le contrôle du risque global par la correction de Bonferroni est adapté à un
nombre restreint de stratégies déclarées à l'avance. Sur un univers de milliers
de variantes essayées, le contrôle de la fouille de données exigerait plutôt un
outil dédié comme le \emph{Reality Check} de White~\cite{white2000}.

Une limite plus subtile tient à l'empilement des corrections. Chaque méthode a
été conçue pour une seule source de dépendance, l'effet de sondage et le
rééchantillonnage pour la dépendance intra-titre, chaque porte pour sa propre
dimension. Prises isolément, elles sont exactes. Utilisées ensemble sur des
données où les dépendances coexistent, elles perdent la garantie de leur seuil
nominal, et l'inférence de la correction mensuelle en particulier n'est plus
exactement calibrée.

Le protocole reste cependant fixe et applique toujours les mêmes étapes, ce
qui varie étant la lecture qu'on en fait. Lorsqu'une dépendance se révèle
négligeable sur un canal, le seuil nominal y demeure valide et la conclusion
peut être chiffrée. Ce n'est que si les trois canaux se manifestaient
fortement en même temps, leurs corrections se cumulant, que le risque global
cesserait d'être chiffrable, et ce cas serait alors signalé. Comme la
dépendance simultanée aux trois canaux reste faible dans nos données, nous
retenons les verdicts obtenus.

Sur le risque lui-même, une précision s'impose. La valeur exacte du risque de
retenir à tort une stratégie n'est pas identifiable, puisque les tests portent
sur les mêmes données et sont donc corrélés. Elle peut en revanche être
majorée, car exiger le franchissement de tous les étages ne peut que réduire
ce risque, lequel reste donc inférieur au seuil d'un étage isolé. Le risque
inverse, celui d'écarter une bonne stratégie, croît au contraire avec le
nombre d'étages et ne se majore pas analytiquement. C'est le témoin qui en
fournit le contrôle, en franchissant l'ensemble du parcours.

Il faut enfin rappeler que la puissance élevée des grands échantillons rend le
moindre écart significatif, si bien que significativité et ampleur de l'effet
doivent être soigneusement distinguées.

Face à chacune de ces limites, le protocole ne tranche jamais en force. Il
préfère écarter une stratégie en signalant la réserve plutôt que de conclure
sur une base fragile.

\section{Perspectives}
\label{sec:perspectives}


La première perspective est la création de nos propres stratégies. Elle
suppose de séparer une période d'apprentissage et une période de test, en
avançant pas à pas dans le temps plutôt qu'en recourant à la validation
croisée classique, laquelle mélangerait le futur et le passé. Le protocole
devra être adapté en conséquence.

Il faudra aussi lui ajouter un garde-fou portant sur les effectifs, qui
écarterait en amont les stratégies trop peu fournies pour que le théorème
central limite s'applique. Ni le test de normalité ni ce théorème ne
permettent alors de conclure, et le verdict doit être tenu pour
ininterprétable plutôt que défavorable.

Une autre voie consiste à modéliser l'avantage en fonction du contexte, par
exemple la probabilité de gain selon le régime de marché ou le secteur
d'activité, ce qui prolongerait directement la piste des stratégies ciblées
évoquée au chapitre~\ref{chap:juger}. Cette voie a été écartée du présent
travail par souci de rigueur, car sur des dizaines de milliers d'observations
dépendantes, les p-values des coefficients d'une régression seraient trop
optimistes, ce qui reproduirait exactement le piège dont ce mémoire traite.
Les outils adaptés, modèles à effets mixtes ou équations d'estimation
généralisées, dépassent le cadre de ce travail.

Le double regroupement de Petersen, évoqué au chapitre~\ref{chap:juger},
relève de la même ambition. Il traiterait d'un seul geste les deux dimensions
que nous avons séparées, ce qui donnerait une version unifiée et plus outillée
du protocole. Nous avons préféré la décomposition, qui rend visible lequel des
canaux fait basculer un verdict, mais l'unification reste l'étape suivante
naturelle.

L'objectif est enfin de fournir à l'agent un maximum d'information, y compris
des signaux tournés vers l'avenir comme les projets de loi ou les signatures
de contrats, ainsi que les relations d'avance et de retard entre titres, au
sens de la causalité de Granger~\cite{granger1969}. Des travaux exploratoires
ont été engagés sur ces fronts. Ce mémoire en constitue la fondation, puisqu'il
fournit l'outil qui permettra de vérifier, à mesure, les stratégies à venir.